\scanpage{33}
\uline{Formulaire de} $L,\Lambda$.

$X$ variété projective lisse connexe polarisée par $\xi_X\in H^2(X)$, $L_X$ opération de cup-produit associée, $\Lambda_X$ opération de degré $-2$ correspondante, $p_i^X$ les projections de $H^*(X)$ sur les morceaux primitifs $P^i(X)$ pour $0\leq i\leq2n$, $Y$ une section hyperplane lisse de $X$, $\xi_Y\in H^2(Y)$ la self-intersection (induite par $\xi_X$), d'où $L_Y$, $\Lambda_Y$, $p_Y^i$; $\phi:Y\to X$ l'immersion, d'où $\phi^*$ et $\phi_*$.

A) On a
\[
\begin{aligned}
1)\quad&\operatorname{Im}L_X=\bigcap_{0\leq i\leq n}\operatorname{Ker}p_i^X,\qquad
\operatorname{Ker}L_X=\sum_{n\leq i\leq2n}\operatorname{Im}p_i^X,\\
2)\quad&\operatorname{Im}\Lambda_X=\bigcap_{n\leq i\leq2n}\operatorname{Ker}p_i^X\qquad
\operatorname{Ker}\Lambda_X=\sum_{0\leq i\leq n}\operatorname{Im}p_i^X,
\end{aligned}
\]
d'où en particulier les formules
\[
\begin{aligned}
1')\quad&p_i^XL_X=0\text{ si }0\leq i\leq n,\qquad L_Xp_i^X=0\text{ si }n\leq i\leq2n,\\
2')\quad&p_i^X\Lambda_X=0\text{ si }n\leq i\leq2n,\qquad\Lambda_Xp_i^X=0\text{ si }0\leq i\leq n,
\end{aligned}
\]
et en particulier
\[
3)\quad p_n^XL_X=L_Xp_n^X=p_n^X\Lambda_X=\Lambda_Xp_n^X=0.
\]
B) On a
\[
4)\quad\Lambda_XL_X=\mathrm{id}_X-\sum_{\deleted{\unsure{0\leq i\leq n}}\inserted{n\leq i\leq2n}}p_i^X,
\qquad L_X\Lambda_X=\mathrm{id}_X-\sum_{\deleted{\unsure{n\leq i\leq2n}}\inserted{0\leq i\leq n}}p_i^X,
\]
\deleted{d'où compte tenu de 1') et 2') les formules
\[
\begin{aligned}
4')\quad&\Lambda_XL_X^2=\deleted{\unsure{L_X}}\inserted{\left(\mathrm{id}_X-\sum_{\unsure{0\leq i\leq2n}}p_i^X\right)L_X},
\quad L_X^2\Lambda_X=\deleted{\unsure{L_X}}\inserted{L_X\left(\mathrm{id}_X-\sum_{\unsure{0\leq i\leq2n}}p_i^X\right)},\\
4'')\quad&L_X\Lambda_XL_X=\unsure{L_X\left(\mathrm{id}_X-\sum_{0\leq i\leq n}p_i^X\right)
=\left(\mathrm{id}_X-\sum_{n\leq i\leq2n}p_i^X\right)L_X}
\end{aligned}
\]
\inserted{et formules ``transposées''.}}

C) On a
\[
5)\quad\phi_*\Lambda_Y\phi^*=\mathrm{id}_X-\sum_{0\leq i\leq2n}p_i^X
\inserted{\underset{\text{(par 4)}}{=}(\Lambda_XL_X+L_X\Lambda_X)-\mathrm{id}_X+p_n^X}
\]
d'où\inserted{,} \deleted{compte tenu de 4''), 1') et 2')
\[
5')\quad\unsure{\frac12|\phi_*}\Lambda_X L_X^2=\unsure{L_X\phi_*\Lambda_Y\phi^*}
\inserted{\deleted{\phi_*\Lambda_Y\phi^*L_X}}
\qquad\unsure{L_X^2\Lambda_X=L_X\phi_*\Lambda_Y\phi^*},
\]}
\inserted{posant
\[
(6)\quad T_X=\mathrm{id}_X+\phi_*\Lambda_Y\phi^*,
\]
(donc $\Lambda_XL_X+L_X\Lambda_X=T_X-p_n^X$),}
d'où compte tenu de \deleted{\unsure{4''}}\inserted{\unsure{3*|3'}}) la formule
\[
\deleted{\unsure{6}}\inserted{7})\quad L_X^2\Lambda_X+2L_X\Lambda_XL_X+\Lambda_XL_X^2
=\deleted{\unsure{4L_X(\mathrm{id}_X-\phi_*\Lambda_Y\phi^*)+2(\mathrm{id}_X-\phi_*\Lambda_Y\phi^*)L_X}}
\inserted{\deleted{\unsure{2L_X+\mathrm{id}L_X}}L_XT_X+T_XL_X},
\]
\inserted{avec $T_X$ \deleted{$\unsure{T_X=\mathrm{id}_X+\phi_*\Lambda_Y\phi^*}$} donné par (6)}

\uline{Corollaire} $(A(X\times X)$ et $C(Y))$ implique $C(X)$.

\margin{On y peut remplacer $A(X\times X)$ par $A^\circ(X\times X)$\par
(condition ``du degré critique'' $H^{\unsure{(2n)-2|(2n)+2}}\longrightarrow H^{(2n)+2}$}
\mspage{32}
